\section{Representation Convergence Hypothesis}
\label{sec:representation_convergence_hypothesis}

The \gls{car} method relies on a assumption, referred to as the \RepresentationConvergenceHypothesisName.
This hypothesis follows the two assumptions established previously throughout this manuscript.

The \MechanisticHypothesisName rejects the notion that deep \gls{nn} constitute irreducible black boxes.
Adopting a mechanistic perspective, we consider that every process results from a sequence of causal computations.
Although these computations may be complex, they remain explainable through the analysis of the internal mechanisms.

The \ConceptHypothesisName states that \gls{nn} achieve generalization by learning concepts.
During optimization, the network progressively organizes its \gls{ls} into representations that capture the information required to solve the task.

Following these two hypotheses, it becomes possible to explain a \gls{nn} by studying the organization of its \gls{ls}.
This family of approaches is referred to as mechanistic interpretability.
Since the \gls{car} method is grounded in these two hypotheses, it falls within this field of research.
More specifically, \gls{car} focuses on a specific aspect of mechanistic interpretability known as universality.
Universality refers to the consistency with which a concept is recovered across different models trained on the same task \cite{templeton2024scaling}.

To provide a explanation for this phenomenon, we developed the \RepresentationConvergenceHypothesisName.
This hypothesis states that, when \gls{nn} are trained to solve the same task, a set of concepts should systematically emerge because these concepts provide a effective way of representing the problem.
As a result, independently trained models should develop similar concepts in their learned representations.

\begin{hypothesis}{Representation Convergence Hypothesis}{convergence}
For a given task, independently trained \gls{nn} tend to develop similar sets of concepts because these concepts provide an efficient representation for solving the task.
\end{hypothesis}

The \gls{car} method is designed to investigate the \RepresentationConvergenceHypothesisName by assessing whether it can consistently recover similar concepts across independently trained models.
The hypothesis therefore serves both as the foundation of the method and as the phenomenon it aims to investigate.
If \gls{car} consistently recovers the same set of concepts across different training runs, this set is more likely to reflect invariant properties of the task than model-specific artifacts.
Conversely, the absence of representation convergence would provide evidence against the hypothesis and would also question the validity of the \gls{car} approach.

It is important to clarify that this hypothesis does not imply convergence at the parameter level.
Independently trained models are expected to converge toward different sets of weights.
The proposed hypothesis concerns a abstract level of representation.
We assume that the conceptual organization induced by their \gls{lr} converges toward an equivalent \gls{cs}, even though the numerical implementation of the networks differ.

The remainder of this chapter develops the methodology used to investigate this hypothesis.
The \ClusteringUniversalityName metric, introduced in Subsection~\ref{subsec:clustering_universality}, provides a quantitative measure of the convergence observed between the concepts identified in independently trained models.